\section{Conclusion}
\label{sec:conclusion}

We described the Semantic Fact Database, a prototype context-indexed
storage system for semantic facts. The system stores each fact atomically
as a section over a finite context poset and answers entity-anchored
queries through a stalk index that returns complete facts in a single hash
lookup. A companion reified triple-store baseline is provided over the
same canonical schema, and a constructive semantic-equivalence proof
binds the two engines together, checked at runtime by an automated harness:
every insert and every query is checked,
online, for canonical result equality.

The evaluation, on five query classes and insert throughput at
\benchscalecount{} scales up to $10^5$ facts, yields a decomposed result
and a clear tradeoff. Against the write-through KG baseline, SFDB is
faster on every query class measured at every scale: entity-anchored
LOOKUP by \rangelookupkg{}, unrestricted GLOBAL
queries by \rangeglobalkg{}, CONTEXT queries --- the class the design
is named for --- by \rangecontextkg{}, and the two temporal-interval
classes by
\rangetemporalkg{} and \rangetemporalunbkg{}. Insert throughput does not
share that pattern: SFDB is slower than KG at the smallest scale
measured and faster at every larger one. The storage-layer control
splits the query-class headline in two: the LOOKUP, GLOBAL, and CONTEXT
advantages survive
against the dict-indexed KG-mem control (\rangelookupkgmem{},
\rangeglobalkgmem{}, and \rangecontextkgmem{}) and are properties of
storing facts atomically
under a context index, while the temporal margins do not survive at
all --- SFDB is at parity or measurably slower than KG-mem on bounded
TEMPORAL at three of four scales --- and are storage-layer effects,
which we report as a negative result rather than reframe. Four
external reference points anchor the comparison outside our own code:
rdflib, on every query class with result
counts verified, and --- the single highest-leverage gap this project's
own reviewer simulation and roadmap identified --- three real, disk-backed
production stores, Apache Jena TDB2, OpenLink Virtuoso, and Ontotext
GraphDB, independently
engineered by three different organisations and accessed over SPARQL/HTTP
exactly as an external client would (Section~\ref{sec:eval:jena},
Section~\ref{sec:eval:virtuoso}, Section~\ref{sec:eval:graphdb}). Against all three, SFDB
remains faster on LOOKUP, GLOBAL, and CONTEXT at every scale, and is overtaken on
TEMPORAL from $10^4$ facts onward: a genuine, scale-dependent
crossover, not a clean win or a hedge, reproduced independently against
three unrelated optimisers, and the first time this project has
measured itself against production engines rather than only its own
research prototypes. Two further comparisons push this further still,
and not only in directions a tidier narrative would have chosen. A
real-world case study
(Section~\ref{sec:eval:wikidata}) against \wikidatanumfacts{} genuine
Wikidata facts and the LUBM standard benchmark workload
(Section~\ref{sec:eval:lubm}) up to \lubmlargescale{} facts both
reproduce the LOOKUP/GLOBAL/CONTEXT
finding cleanly, while TEMPORAL lands somewhere different each time ---
beating both in-house baselines on Wikidata, roughly at parity with
them on LUBM's largest scales --- evidence that TEMPORAL alone among
the four query classes is genuinely sensitive to a dataset's actual
shape. Building the LUBM comparison surfaced something larger than a
data-shape finding: a genuine performance bug in SFDB's own CONTEXT
query resolution, present for the entire evaluation and exposed only by
a topology shape (many large, mutually unrelated top-level contexts)
none of this paper's other datasets have, which required re-running and
re-verifying every CONTEXT number this paper reports and retracting two
earlier claims built on the unfixed numbers. We report the retraction
alongside the fix rather than quietly substitute corrected numbers and
leave the earlier narrative standing. The
tradeoff that funds the surviving advantages is memory: SFDB uses
\rangememkg{} more resident memory per fact than the reified-triple
baseline, and \rangememkgmem{} more than the control. We report this
plainly. We also report --- once, in
full, in Section~\ref{sec:discussion:threats} --- that ten bugs,
performance gaps, and measurement findings were surfaced while preparing this evaluation,
by eight different mechanisms, only one of which was the automated
verification harness catching itself. We treat the fact that our own
evaluation needed ten rounds of correction --- one of them large enough
to change the paper's own central crossover claim from two query
classes to one --- as part of the paper's
evidence for the value of that scrutiny rather than as an embarrassment
to suppress.

The formal apparatus in the paper is that of presheaves on a finite poset.
We do not claim a novel sheaf-theoretic theorem: on the Alexandrov topology
that fits our setting, the sheaf condition holds vacuously and could not
serve as a technical contribution. The value of the framework is
operational. It gives us a disciplined vocabulary for restriction, stalks,
and global sections that maps directly onto the components of the
implementation, and it gives us a canonical intermediate representation
that makes cross-engine equivalence a checkable property rather than an
aspiration.

We do not claim that the design supplants production RDF stores, and the
Jena TDB2, Virtuoso, and GraphDB comparisons are direct evidence for
why not: SFDB
loses to all three on one of the five query classes measured once the
working set is large enough. Our KG
baseline remains a controlled research implementation, and the
LOOKUP/GLOBAL/CONTEXT speedups against it in particular
are best interpreted as a measurement of what disciplined context-indexed
storage buys over a strict reification-plus-joins design at comparable
engineering effort, now corroborated by three independent production-store
results and the LUBM standard workload rather than merely asserted. Extending the evaluation to
further
standard benchmark workloads (LUBM at the much larger scales published
LUBM results use, BSBM, WatDiv) remains future work, to check
whether the crossover generalises further still. The artifact and reproducibility
manifest for every result in the paper are released with the paper.
